\documentclass[10pt]{article}
\usepackage[T1]{fontenc}
\usepackage{lmodern,amsmath,amssymb,amsthm,microtype}
\usepackage[margin=0.84in]{geometry}
\usepackage[colorlinks=true,urlcolor=blue,citecolor=blue]{hyperref}
\newtheorem{theorem}{Theorem}
\title{The Greedy Digital-Root Sequence Is a Permutation}
\author{Qichao Wang\\{\small Hebei University of Technology}\\[0.3em]
Daoyu Dong\\{\small University of Electronic Science and Technology of China}}
\date{October 26, 2026}
\hypersetup{pdftitle={A Permutation Proof for OEIS A340138},pdfauthor={Qichao Wang and Daoyu Dong}}
\begin{document}
\maketitle
\begin{abstract}
We prove the permutation conjecture for OEIS A340138. The argument
uses the finite set of digital roots: any root occurring infinitely
often forces every other root to do so, and then every positive
integer must appear. The same proof applies in every integer base
at least two. We explicitly resolve an indexing inconsistency in
the entry by following its published terms and Mathematica program.
\end{abstract}

\section{The precise sequence}
For an integer base $b\ge2$, let $\operatorname{dr}_b(0)=0$, and
for $x>0$ let
\[
 \operatorname{dr}_b(x)=1+((x-1)\bmod(b-1))\in\{1,\ldots,b-1\}.
\]
This is the usual iterated digit sum. Start with $a_b(1)=0$ and,
for $n\ge1$, take $a_b(n+1)$ to be the least unused nonnegative
integer containing the digit $\operatorname{dr}_b(a_b(n))$ in its
ordinary base-$b$ expansion. Thus $a_b(2)=b$.

At $b=10$ these are exactly the displayed terms and the Mathematica
program in OEIS A340138 \cite{oeis}. The entry's prose instead
puts $a(n-1)$ in the rule for $a(n+1)$ while specifying only $a(1)$;
that rule is not initialized and disagrees with its program. We use
the immediate-predecessor indexing above, as established by both
the program and data. The permutation assertion below concerns
that published sequence; it does not silently supply an unspecified
initial value for a different two-step rule.

\begin{theorem}
For every $b\ge2$, the sequence $(a_b(n))_{n\ge1}$ is a permutation
of the nonnegative integers. In particular the permutation
conjecture in A340138 is true for its published sequence.
\end{theorem}
\begin{proof}
Every stage is well-defined: infinitely many base-$b$ integers
contain any specified digit, while only finitely many have been
used. Terms are distinct by construction, and every term after
the initial zero is positive.

For $1\le r<b$, let $\mathcal C_r$ be the positive integers
containing the digit $r$. If the digital root $r$ occurs at
infinitely many terms, then every $M\in\mathcal C_r$ must appear.
Indeed, if $M$ stayed unused, every step immediately following
a term of root $r$ would choose an unused member at most $M$.
The infinitely many such choices are distinct, contradicting
the finiteness of $\{1,\ldots,M\}$.

At least one positive root $r$ occurs infinitely often, since
the set of possible positive roots is finite. Fix any other
root $s\in\{1,\ldots,b-1\}$, and choose the unique
$d\in\{0,\ldots,b-2\}$ satisfying
$r+d\equiv s\pmod{b-1}$. For every $L\ge1$, the base-$b$ word
\[
 r\,\underbrace{0\cdots0}_{L\text{ zeros}}\,d
\]
represents a distinct positive member of $\mathcal C_r$ with
digital root $s$. The preceding paragraph forces all these
integers to appear. Therefore $s$ also occurs infinitely often.
This proves that every positive root recurs infinitely often.

Finally every positive integer contains at least one nonzero
digit $s$, so it belongs to $\mathcal C_s$ and must appear.
The initial zero already appears exactly once. This proves the
permutation assertion.
\end{proof}

\section{Independent verification and scope}
The accompanying \texttt{verify.py} compares all 10,000 published
decimal terms with an exact implementation. A separate direct
candidate scan, using repeated digit sums rather than the modular
digital-root formula, verifies the first 200 terms in bases 2
through 16. The constructive root-propagation words are checked
for every pair of roots in those bases and lengths 1 through 20.
These computations validate the implementation and the intended
indexing; the proof establishes the entire infinite permutation.
No finite computation is used to infer surjectivity.

\begin{thebibliography}{1}
\bibitem{oeis} M. De Vlieger, OEIS A340138,
\url{https://oeis.org/A340138}, accessed October 26, 2026.
\end{thebibliography}
\end{document}
